\documentclass[11pt,a4paper]{article}
\usepackage{cmap}
\usepackage[utf8]{inputenc}
\usepackage[T1]{fontenc}
\usepackage{lmodern}
\usepackage{microtype}
\usepackage[margin=1.0in]{geometry}
\usepackage{amsmath,amssymb,amsthm,mathtools}
\usepackage{booktabs,array,tabularx}
\usepackage[svgnames]{xcolor}
\usepackage{listings}
\usepackage[obeyspaces,spaces]{url}
\usepackage[colorlinks=true,linkcolor=DarkBlue,citecolor=DarkBlue,urlcolor=DarkBlue]{hyperref}

\urlstyle{tt}
\def\UrlFont{\ttfamily\upshape}
\renewcommand{\arraystretch}{1.3}
\setlength{\tabcolsep}{5pt}

\sloppy
\setlength{\emergencystretch}{4em}
\hfuzz=4.0pt
\vfuzz=4.0pt

\newtheorem{theorem}{Theorem}[section]
\newtheorem{lemma}[theorem]{Lemma}
\newtheorem{proposition}[theorem]{Proposition}
\newtheorem{corollary}[theorem]{Corollary}
\theoremstyle{definition}
\newtheorem{definition}[theorem]{Definition}
\newtheorem{model}[theorem]{Modelling assumption}
\theoremstyle{remark}
\newtheorem{remark}{Remark}[section]

\lstdefinestyle{lean}{
  basicstyle=\ttfamily\footnotesize, breaklines=true, frame=single,
  framesep=4pt, xleftmargin=4pt,
  literate={≤}{{$\leq$}}1 {∀}{{$\forall$}}1 {→}{{$\to$}}1 {↔}{{$\leftrightarrow$}}1
           {ℚ}{{$\mathbb{Q}$}}1 {α}{{$\alpha$}}1 {¬}{{$\neg$}}1 {·}{{$\cdot$}}1,
}
\lstset{style=lean}

\title{\textbf{Proof-Carrying Detection of Reticulate Evolution:\\
a Tree-Level Null, the Measured Cost of the Alignment Step,\\
and the Gromov Hinge of Buneman's Converse}}
\author{Xavier Callens\\
\small SocrateAI Scientific Research \& Formal Verification Laboratory\\
\small Formal artifacts: \texttt{ANSE.GeneticaReticulata},
\texttt{ANSE.GeneticaReticulataTree} (Lean 4)}
\date{Edition 4, 23 October 2026. Closes the topological line on measured
grounds (\S\ref{sec:topology}), measures the maximum statistic's defect across
taxon counts rather than arguing it from two points (\S\ref{sec:scale}), and
reports a filtration funnel in which no variant met its pre-registered
criterion. \textbf{Not peer reviewed.}}

\begin{document}
\maketitle

\begin{abstract}
\noindent
Tests for recombination are normally calibrated against a permutation ensemble.
We develop one calibrated against a \emph{theorem}, report the two ways our own
first version of it was wrong, and give the measurement that we believe is new.

A strict molecular clock on a clonal genealogy makes expected $p$-distances
ultrametric, and an ultrametric admits no unfilled Vietoris--Rips cycle at any
scale; three sequences violating the strong triangle inequality are therefore a
\emph{certificate}, not a surprise relative to shuffles. We formalise the
inference rule in Lean~4 and emit certificates that embed the aligned rows, so
the kernel recomputes every count by \texttt{rfl} rather than trusting a search.

The statistic this licenses is nonetheless the wrong one, and a control we had
not run says so. On a clonal tree with a \textbf{relaxed clock and no
reticulation whatever}, divergence-matched so that only the clock differs, the
ultrametric margin exceeds every single-recombinant margin in the positive
control; at 32 leaves $\min(\text{no reticulation}) = 99/4000 > 79/4000 =
\max(\text{one recombinant})$ at every seed. The four-point condition encodes
\emph{tree} without \emph{clock}, we prove in Lean that every symmetric
ultrametric satisfies it, and the quartet margin on the same data does not move.
So the quartet statistic must carry the method, and the licence for reading the
gap between the two conditions is a theorem rather than a convention.

Running this on real data required writing the alignment step, which puts an
aligner on the trust surface that no downstream proof covers. We measure it
instead of caveating it: on simulated populations with no reticulation and a
known true alignment, reference-anchored alignment manufactures $0$ of $9100$
four-point violations at a $0.4\%$ gap load and $1600$ of $9100$ at $8.2\%$. We
find no published benchmark that varies alignment quality as an independent
variable, so we believe this is the first direct measurement of a mechanism the
field already asserts. Read as an admissibility threshold it refuses $7$ of $10$
real viral populations drawn from a public corpus; the $3$ that pass agree with a
literature prior fixed in advance, and the two certificates clearing their matched
floor are accepted by the kernel with both corrupted twins rejected.

Two corrections to our own apparatus are reported as results. The band-edge test
by which we had certified the aligner is not a soundness certificate, and it had
really failed: a pair scoring $16555$ against an unbanded optimum of $17275$ with
the band edge never touched. Certifying every alignment against the unbanded
optimum dropped the figure above from $4876$ to $1600$ --- two thirds of our
headline was an uncertified step inside the measuring apparatus. Separately, the
statistic we report is a maximum over subsets where the established treelikeness
tests are averages; run seed-paired against the hard null, the maximum separates
$16.57\times$ against $5.07\times$ for mean $\delta_q$ and $1.22\times$ for the
proportion of violating quadruples.

Edition~1 listed the missing head-to-head as its chief limitation. We now run it
against the one incumbent whose null also assumes no clock, the pairwise
homoplasy index, on the same alignments under the same deletion policy, and the
result revises what this method may claim. \textbf{On the five populations where
both tests can be measured the two agree exactly}: both are specific on the
strict-clock \emph{and} the relaxed-clock null, and both detect all three
reticulate populations at every seed. The incumbent is not beaten on power, so
the case for a certificate cannot rest on power. It rests instead on two things
the comparison does not touch --- an exhibitable witness a kernel recomputes, and
the absence of a permutation null, which matters because the artefact class we
measure in Section~\ref{sec:aligner} is reported to be invisible to exactly that
kind of null.

Finally we prove the hinge of Buneman's converse, which this programme had only
cited: a symmetric dissimilarity satisfies the four-point condition if and only
if its Gromov product satisfies the max inequality at every root. We find no
formalisation of tree metrics, splits or Buneman's theorem in Lean, Isabelle or
Coq, so we believe this is first of its kind. All Lean declarations compile with
zero \texttt{sorry} and no axioms beyond
$\{\texttt{propext}, \texttt{Classical.choice}, \texttt{Quot.sound}\}$.
\end{abstract}

\tableofcontents

\section{Introduction}

A test for recombination has to answer the question: compared with what? The
standard answer is an ensemble --- shuffle the columns of the alignment, recompute
the statistic, and read off a $p$-value. It is a serviceable answer with two
costs. The first is a resolution floor: with $N$ shuffles no $p$-value below
$1/(N+1)$ is attainable, so the apparatus has an opinion about what counts as
significant before the data arrives. The second is sharper and less often
stated. The manual of the field's dominant tool observes that misalignment
produces false recombination signals which \emph{``will be detected as
recombination events with significant p-values''} while \emph{``these false
positive signals will likely be undetectable in the shuffled
alignments''}~\cite{rdp5}. Column shuffling destroys the spatial clustering of
misalignment along with the signal, so an entire class of artefact survives the
real analysis and vanishes from the null. A permutation null cannot see it.

This paper develops the alternative: a null that is a theorem, so that a
violation is a witness rather than a tail probability.

\begin{model}[What the clock-level null assumes]\label{mod:clock}
Write $d(x,y)$ for the $p$-distance between two aligned sequences --- differing
columns over columns comparable in both. Expected $p$-distances satisfy
\begin{equation}\label{eq:ultra}
  d(x,z) \;\le\; \max\!\big(d(x,y),\, d(y,z)\big) \qquad \text{for all } x,y,z
\end{equation}
exactly when descent is \textbf{clonal} \emph{and} the molecular clock is
\textbf{strict}. Two assumptions, one inequality.
\end{model}

If \eqref{eq:ultra} holds then $d(\cdot,\cdot) \le \epsilon$ is transitive, so
every connected component of the Vietoris--Rips complex at scale $\epsilon$ is a
clique and no unfilled cycle exists at any scale. Three sequences with
$\max(d(x,y), d(y,z)) < d(x,z)$ therefore contradict a universally quantified
statement. Nothing is calibrated.

Section~\ref{sec:nulls} states the two nulls and proves they nest.
Section~\ref{sec:refute} refutes the clock-level statistic with a control that
contains no reticulation. Section~\ref{sec:deletion} deals with missing data.
Section~\ref{sec:aligner} is the paper's main empirical contribution: the
alignment step is on the trust surface, and this is how large it is.
Section~\ref{sec:band} reports the proxy inside our own apparatus that failed and
what it had inflated. Section~\ref{sec:maxmean} prices the maximum against the
average, and Section~\ref{sec:headtohead} runs the comparison against the
incumbent that edition~1 was missing. Section~\ref{sec:real} runs the test on
real viral populations.
Section~\ref{sec:gromov} proves the hinge of Buneman's converse.
Section~\ref{sec:limits} says what would make all of it wrong.

\paragraph{What is and is not claimed.} We claim: the nesting theorem and the
Gromov equivalence, both machine-checked; the refutation of the clock-level
statistic on simulated data; the measurement of alignment-manufactured
violations; and three real-data verdicts that agree with priors fixed in advance.
We do not claim a validated method. Three admissible populations is not a
validation, no head-to-head against the established tests has been run, and
Buneman's converse remains cited for its second half. This paper has not been
peer reviewed.

\section{Two nulls, and the fact that they nest}\label{sec:nulls}

\begin{definition}[Four-point condition]\label{def:fp}
A symmetric $d$ satisfies the \emph{four-point condition} when for all
$a,b,c,e$, not necessarily distinct,
\begin{equation}\label{eq:fp}
  d(a,b) + d(c,e) \;\le\;
  \max\!\big(d(a,c)+d(b,e),\; d(a,e)+d(b,c)\big).
\end{equation}
\end{definition}

A quadruple admits three pairings; \eqref{eq:fp} says no pairing's distance sum
strictly exceeds both others. This is a statement about tree structure with no
constraint on rates: branch lengths are free.

\begin{theorem}[Tree-Metric Theorem; Buneman. \emph{Cited, not proved here}]\label{thm:buneman}
Let $\delta$ be a nonnegative dissimilarity map on a finite set --- zero on the
diagonal, symmetric, non-negative. Then $\delta$ is a tree metric if and only if
it satisfies \eqref{eq:fp}~\cite{buneman,semplesteel}.
\end{theorem}

The hypotheses are weaker than ``metric'', and the distinction matters enough to
state: setting $y = z$ in \eqref{eq:fp} yields the triangle inequality, so the
triangle inequality is a \emph{consequence} of the condition and requiring it as
a hypothesis would state something strictly weaker than the theorem. An earlier
draft of our own documentation made that error.

\begin{theorem}[The nulls nest]\label{thm:nest}
Every symmetric ultrametric satisfies the four-point condition.
\end{theorem}

\begin{proof}
Fix $a,b,c,e$. Four applications of \eqref{eq:ultra}, with symmetry to order each
application, give $d(a,b) \le \max(d(a,c), d(b,c))$,
$d(a,b) \le \max(d(a,e), d(b,e))$, $d(c,e) \le \max(d(a,c), d(a,e))$ and
$d(c,e) \le \max(d(b,c), d(b,e))$. Each is a disjunction, giving $2^4$ cases;
splitting further on the order of the four bounding quantities, every branch
bounds $d(a,b)+d(c,e)$ by one of the two right-hand sums of \eqref{eq:fp} by a
linear arithmetic step. The case analysis is mechanical and is discharged by the
kernel (Listing~\ref{lst:nest}).
\end{proof}

\begin{corollary}[What a quartet violation buys]\label{cor:q}
If $d$ is symmetric and fails \eqref{eq:fp} then $d$ is not ultrametric either.
\end{corollary}

Corollary~\ref{cor:q} is the licence for the reading that organises
Section~\ref{sec:refute}: large triple margins together with floor-level quartet
margins indicate rate variation on a tree, because the nesting leaves that the
only region available.

\begin{lstlisting}[caption={\texttt{formal/ANSE/GeneticaReticulataTree.lean}. The
\texttt{rcases} chain is the $2^4$ case split, with four further splits on the
order of the bounding quantities.},label={lst:nest}]
def IsFourPoint (d : α → α → ℚ) : Prop :=
  ∀ a b c e, d a b + d c e ≤ max (d a c + d b e) (d a e + d b c)

theorem fourPoint_of_ultrametric {d : α → α → ℚ}
    (hsym : ∀ x y, d x y = d y x) (hu : IsUltrametric d) : IsFourPoint d
\end{lstlisting}

The module carries 18 declarations through an axiom gate that compiles each and
runs \texttt{\#print axioms}: zero \texttt{sorry}, and no axioms beyond
$\{\texttt{propext}, \texttt{Classical.choice}, \texttt{Quot.sound}\}$. Checking
compilation alone is insufficient --- both \texttt{sorry} and a locally declared
axiom exit zero --- which is why the gate inspects the axiom list.

\section{The clock-level statistic is refuted by a control with no reticulation}
\label{sec:refute}

Condition \eqref{eq:ultra} needs three sequences rather than four, so it is
cheaper, and our first version reported it. It separated a simulated clonal
population from a simulated recombinant cleanly. The separation was real and the
inference from it was wrong, for the reason
Modelling~assumption~\ref{mod:clock} warns of: a violation tells you one of two
assumptions failed, not which.

The decisive control is a \textbf{divergence-matched relaxed clock}: clonal
descent on a balanced tree, no reticulation whatever, with each lineage's
substitution rate multiplied by $e^{su}$ for $u$ uniform on $[-1,1]$. The
matching is not optional. Since $\mathbb{E}[e^{su}] = \sinh(s)/s > 1$, an
unnormalised relaxed clock is also a \emph{faster} clock, and a violation could
then be attributed to saturation rather than to rate variation; dividing by that
mean leaves expected divergence equal to the strict-clock case while realised
rates vary widely. In our first run the normalisation changed the answer, taking
the relaxed-clock quartet violation count from $2/5$ seeds to $0/5$.

\begin{table}[t]
\centering\small
\caption{Thirty-two leaves, $4000$ nt, three seeds, exhaustive over all
$35\,960$ quadruples per dataset. Margins are exact rationals with denominator
$4000$; numerators given seed by seed.}
\label{tab:refute}
\begin{tabular}{@{}lccll@{}}
\toprule
simulated population & reticulation? & clock & triple margin & quartet margin \\
\midrule
clonal                     & no  & strict  & $24, 30, 29$   & $10, 9, 10$ \\
clonal                     & \textbf{no}  & \textbf{relaxed} & $\mathbf{155, 157, 99}$ & $11, 18, 8$ \\
one recombinant            & yes & strict  & $79, 77, 75$   & $79, 70, 69$ \\
two reciprocal mosaics     & yes & strict  & $79, 77, 75$   & $\mathbf{158, 154, 126}$ \\
\bottomrule
\end{tabular}
\end{table}

Read rows two and three of Table~\ref{tab:refute} together. The smallest triple
margin produced by a population with \emph{no reticulation in it at all} is
$99/4000$; the largest produced by a population containing a real recombinant is
$79/4000$. The ordering is inverted at every seed. On this data a larger triple
margin is weak evidence \emph{against} reticulation. The quartet column does not
move: the relaxed clock stays at the strict clock's level while the
reciprocal-mosaic positive control separates from it by a factor of $12.6$ at its
weakest seed, with no overlap. At sixteen leaves the same experiment gives a
relaxed-clock triple margin reaching $106/4000$ --- above every
single-recombinant margin --- exceeding twice the clonal floor at $3$ of $5$
seeds, while the quartet statistic does so at $0$ of $5$.

Adding taxa makes the triple statistic worse, as the mechanism predicts: it is a
maximum over triples, so more leaves give more chances at an extreme pair of
rates.

\begin{remark}
This confound is not peculiar to our statistic. Frankel and An\'e report that
rate variation across lineages produces a marked increase in the false discovery
of reticulation for three summary tests of introgression, with the $D_3$
statistic reaching about $80\%$ type-I error, and read this as $D_3$ responding
more strongly to a departure from the clock than to reticulation
itself~\cite{frankelane}. A null that assumes a clock is vulnerable wherever it
is used. The established test whose null assumes no clock, Bruen, Philippe and
Bryant's $\Phi_w$, reports $5.7\%$ false positives on a non-recombinant control
where NSS reports $37.5\%$~\cite{bruen}.
\end{remark}

\paragraph{Recommendation.} Report the quartet statistic as primary. Keep the
triple statistic --- it is cheaper, and a triple violation with \emph{no} quartet
violation is itself informative, being the signature of rate variation on a tree
--- but a triple certificate alone is not evidence of reticulation.

\section{Missing data: complete deletion within the subset}\label{sec:deletion}

Pairwise deletion computes each pair's distance over the columns comparable in
that pair, giving different pairs different denominators, and that alone can
break \eqref{eq:ultra} on a genealogy that is exactly a clock tree. The fix is
\emph{complete deletion within the subset}: count only columns at which every
member of the triple or quadruple carries a definite base. All distances in a
certificate then share one denominator, so the obligations become comparisons
between integers with no cross-multiplication --- a smaller trust surface and the
repair at once. The cost is power: at a $25\%$ gap load roughly $1000$ to $1700$
of $4000$ columns survive within a triple.

Measured on gap-blocked but strictly clonal, strictly clock-like populations:
pairwise deletion calls $33$ to $52$ of the $560$ triples violations that
complete deletion within the triple denies, at every one of five seeds; at
thirty-two leaves, $286$ to $366$ of $4960$. With no missing data the two
policies agree on all $560$, as they must.

\begin{remark}[Margins under the two policies are not comparable]
Complete deletion shrinks the denominator, so one differing site is worth more
and its margins come out \emph{larger} although the policy is strictly more
conservative. A criterion comparing the two sets of margins was written into our
first run and withdrawn. Compare per-subset \emph{verdicts}, not magnitudes.
\end{remark}

\section{The alignment step is on the trust surface}\label{sec:aligner}

A certificate is machine-checked, which invites the claim that the conclusion
does not depend on the pipeline. The precise statement is narrower. The kernel is
handed the sequences and recomputes every count from them, so a false
inequality, a wrong count, a transposed pair or a miscomputed margin is excluded.
What remains trusted is that the embedded rows are the data they are claimed to
be --- and on real data those rows come out of an \textbf{aligner}. If the
alignment places non-homologous residues in the same column, the counts are
computed faithfully from the wrong data and the kernel accepts them. Proof moves
the trust boundary; it does not close it.

That cannot be proved away, so we measure it. We simulate clonal, relaxed-clock
descent with deletion events in a fixed coordinate frame, so the frame \emph{is}
the true alignment and \textbf{there is no reticulation in the population at
all}; strip the gaps; hand the ragged sequences to the aligner; realign; and
compare per-quadruple \emph{verdicts}, not margins, since the two alignments
have different widths. Any quadruple the inferred alignment calls a violation
and the true alignment does not was manufactured by the aligner.

The aligner is reference-anchored: each sequence is aligned to a centroid
reference by affine-gap global dynamic programming and projected onto the
reference's coordinates. This is the construction viral genomic surveillance
uses, and the final step of MUSCLE5's \texttt{Super5} mode. Its one structural
advantage is that without a guide tree there is no error-propagation failure
mode; its costs are that insertions relative to the reference are dropped and
that every pair of non-reference sequences is related through the reference.

\begin{table}[t]
\centering\small
\caption{The admissibility ladder. Sixteen taxa, $4000$ nt, five seeds, $1820$
quadruples per seed ($9100$ per rung). Every alignment certified optimal
(Section~\ref{sec:band}).}
\label{tab:ladder}
\begin{tabular}{@{}cccrc@{}}
\toprule
mean deletion & inferred gap & worst column & quadruples & seeds \\
length (nt) & fraction & recovery & manufactured & with any \\
\midrule
 30 & $0.00000$ & $1.00000$ & $0$    & $0/5$ \\
 30 & $0.00398$ & $0.99997$ & $0$    & $0/5$ \\
 30 & $0.00549$ & $0.99880$ & $251$  & $2/5$ \\
 30 & $0.00870$ & $0.99712$ & $108$  & $1/5$ \\
 30 & $0.01507$ & $0.99906$ & $110$  & $3/5$ \\
 30 & $0.02012$ & $0.99831$ & $149$  & $3/5$ \\
120 & $0.00000$ & $1.00000$ & $0$    & $0/5$ \\
120 & $0.01322$ & $0.99901$ & $57$   & $2/5$ \\
120 & $0.01627$ & $0.99901$ & $54$   & $2/5$ \\
120 & $0.03731$ & $0.99467$ & $324$  & $3/5$ \\
120 & $0.05254$ & $0.97501$ & $253$  & $3/5$ \\
120 & $0.08233$ & $0.96032$ & $\mathbf{1600}$ & $5/5$ \\
\bottomrule
\end{tabular}
\end{table}

At the heaviest rung the aligner invents $1600$ of the $9100$ quadruple verdicts
on a genealogy that is exactly a tree. \textbf{Past a threshold, alignment error
is not a correction to this test; it is the result.} The curve is not monotone,
which is what a maximum over $1820$ quadruples at five seeds looks like, and is
one reason to gate on \emph{whether} anything was manufactured rather than on how
much.

So the admissibility threshold is read off the ladder rather than chosen: a gap
fraction of $0.00398$ and $8.06 \times 10^{-4}$ dropped insertions per aligned
cell, the heaviest rung at which nothing was manufactured at any seed. Both are
\emph{rates}; an absolute insertion count is not comparable between a $3$~kb and
a $12$~kb genome.

\paragraph{Why we believe this measurement is new.} That misalignment
manufactures recombination signal is not a novel claim and not a contested one.
It is stated outright in the RDP5 manual --- \emph{``misalignment of sequences is
a major cause of false recombination signals''} and \emph{``all recombination
signal detection methods are extremely sensitive to misalignment and will
usually identify misaligned regions of sequence as having a recombinant
origin''} --- and misalignment artefact is a named output category in that
program~\cite{rdp5}. What we did not find is the quantification. Across the
recombination benchmarking literature we examined, every study varies divergence,
rate heterogeneity, topology, selection or sequence count while holding the
alignment fixed; the most recent full benchmark we reached does not discuss gaps
at all~\cite{jaya}. A measurement of the form ``alignment error alone
manufactures $N$ of $M$ apparent signals on data containing no reticulation''
appears therefore to be the first direct measurement of a mechanism the field
already accepts. That is confirmatory-but-novel, and we state it that way rather
than as a discovery. The nearest quantified precedent we found is in an adjacent
task: aligner choice driving a reported $48$--$82\%$ of false positives in
\emph{Drosophila} positive-selection scans.

\paragraph{An honest limitation of the simulation.} The deletion rate and mean
deletion length are \emph{chosen}, not calibrated: we found no published viral
indel length distribution and no clean per-site viral indel rate. The ladder
therefore varies a load; it does not match one to biology.

\section{Certifying the band: a proxy inside our own apparatus}\label{sec:band}

The dynamic programme is banded, because $12$~kb against $12$~kb is
$1.4 \times 10^8$ cells and the optimal path of two homologues does not wander
far from the diagonal. A band is a cap on how much indel a pair may be assigned,
so it is a place where an aligner can silently return a sub-optimal path. Our
first implementation detected when the optimal traceback \emph{touched} a band
edge, widened the band until it did not, and documented this as a guarantee that
no constrained path was ever returned without saying so.

That documentation was wrong. Gibrat measures both failure directions on real
data: $2.2\%$ of alignments are optimal yet reach the band edge, and $0.5\%$ are
sub-optimal without ever reaching it~\cite{gibrat}. Not touching the edge is
evidence of optimality, not proof of it.

The repair needs no theorem. Computing the unbanded optimum by full Gotoh
dynamic programming in $O(nm)$ time and $O(m)$ memory --- no band, no heuristic,
no traceback --- and refusing any banded path that does not attain it gives a
certificate for one extra pass per pair, which at fifteen pairs per population is
affordable. We also note, and do not use, a cheaper sufficient test: Gibrat gives
an upper bound on the score of any path leaving the band, so attaining that bound
certifies optimality in constant time. His derivation is for a band about the main
diagonal while ours is about the proportional diagonal, so the bound does not
transfer unexamined, and recomputation needs no derivation to be trusted.

\paragraph{The proxy had really failed.} The first certified run exited non-zero
immediately, on a pair scoring $16555$ against an unbanded optimum of $17275$ ---
a deficit of $720$ points, with the traceback never touching an edge, so
edge-driven widening never fired. The widening loop is now driven by the
\emph{score}: widen on either signal, and return only a path that attains the
optimum, or an error.

\paragraph{What it had inflated.} With every alignment proved optimal,
Table~\ref{tab:ladder}'s heaviest rung manufactures $1600$ rather than $4876$,
and worst column recovery is $0.960$ rather than $0.868$. Roughly two thirds of
what we had reported as reference-anchored alignment error was our own unchecked
heuristic. The qualitative conclusion is unchanged and the thresholds do not
move, so no verdict in Section~\ref{sec:real} changes; but the magnitude was
overstated threefold, in the direction that makes a finding look more impressive.
We record this because it generalises: \textbf{an uncertified step inside a
measuring apparatus inflates the thing being measured}, and it will tend to
inflate it favourably, which is the direction least likely to prompt a second
look.

\begin{remark}[A textbook algorithm that is wrong in print]
While checking the above we found that Gotoh's \emph{published} initialisation is
incorrect --- it omits one of two gap-opens and can return a score that is not
the cost of any alignment --- and that the error is reproduced in standard
textbooks and in five of ten inspected implementations~\cite{flouri}. Ours turns
out to be the corrected form, but nothing in our codebase established that until
the dynamic programme was tested against exhaustive enumeration of every
alignment of short sequences rather than against another dynamic programme.
\end{remark}

\begin{remark}[A tie-break that is a free parameter]
The control that the aligner is the identity on gapless input failed at one seed
of five on our first implementation, and not because of a bug. With an affine gap
of $-10$ to open and $-1$ per residue, three consecutive mismatches score
$3 \times (-4) + 5 = -7$, and deleting one reference residue, aligning the rest,
then inserting one query residue scores $-11 + 15 - 11 = -7$. The identity was
one of two optimal alignments and the implementation returned the other. The
repair was to \emph{state} the preference. An aligner has free parameters that no
amount of downstream proof constrains.
\end{remark}

\section{A maximum, where the established tests are averages}\label{sec:maxmean}

Our certificate is a maximum over quadruples: the single most-violating
quadruple, exhibited as a witness the kernel can check. That is what makes a
certificate possible, and it is also an extreme-value statistic whose null level
grows with the number of subsets examined. Our own grid shows it --- the quartet
floor moves from $5$--$7$ at sixteen taxa to $9$--$10$ at thirty-two --- so
margins are not comparable across populations of different size.

The treelikeness literature took the other branch: its statistics are averages
over quadruples, among them $\delta_q = (m_1-m_2)/(m_1-m_3)$ for sorted pairing
sums $m_1 \ge m_2 \ge m_3$~\cite{deltaplots}. We therefore ran the comparison.

\begin{table}[h]
\centering\small
\caption{Seed-paired against the hard null (clonal, relaxed clock, no
reticulation). Separation is the weakest seed's ratio of the reciprocal-mosaic
control to that null.}
\label{tab:maxmean}
\begin{tabular}{@{}llcr@{}}
\toprule
statistic & what it is & beats the null & worst-seed separation \\
\midrule
quartet max margin & exact, with a witness & $5/5$ & $\mathbf{16.57\times}$ \\
mean $\delta_q$ & the averaged statistic & $5/5$ & $5.07\times$ \\
proportion violating & share of quadruples & $5/5$ & $1.22\times$ \\
\bottomrule
\end{tabular}
\end{table}

All three rank both reticulate populations above the relaxed-clock null at every
seed, so none of them inherits the confound that destroyed the triple statistic.
At fixed taxon count the maximum separates more than three times as widely as the
mean, and the proportion-based variant barely separates at all. We note that this
reproduces, independently, a published finding that the proportion-of-resolved-%
quadruples variant fails as a test while the averaged variants hold a correct
false-positive rate.

What the comparison does \emph{not} settle is the maximum's actual weakness,
which is incomparability across taxon counts rather than poor separation within
one; our grid is all at one size. Both statistics are therefore reported, and
only the maximum is exact and kernel-checkable, because a mean has no witness to
exhibit. That is the trade, and it is now priced rather than assumed.

\section{The head-to-head against the incumbent}\label{sec:headtohead}

A method should be compared with what it hopes to improve on, and edition~1 of
this paper did not do that. The relevant incumbent is not one of the many tests
that assume a clock --- those fail the control of Section~\ref{sec:refute} by
construction --- but the one whose null assumes \textbf{site-order
exchangeability} and nothing about rates: the pairwise homoplasy index of Bruen,
Philippe and Bryant~\cite{bruen}. Like the four-point condition it is immune by
design to the relaxed-clock confound, so a comparison between them is about
power and about what each licenses, not about survival.

We implement a restricted form. Two biallelic sites are \emph{incompatible} when
all four combinations of their states occur among the sequences --- the
four-gamete test. Under a single tree, compatibility between two sites does not
depend on their separation in the alignment; under recombination it does, since
neighbouring sites tend to share a local history. The statistic $\Phi_w$ is the
mean compatibility of informative sites within a window $w$, and the null comes
from permuting the site order. Two restrictions are stated rather than hidden:
only biallelic parsimony-informative columns are used, so the four-gamete test is
exact and needs no multistate refinement; and columns with a gap in any member
are skipped, which is the same complete-deletion policy the certificate uses, so
both tests see the same data. This is not the published statistic in full and we
do not present it as such.

One property of the implementation is worth noting. $\Phi_w$ is a count of
compatible neighbouring pairs over a count of pairs, and the pair count is
invariant under permutation of the site order, so the permutation test compares
\emph{numerators}: it is exact integer arithmetic, like the certificate, rather
than a normal approximation.

\begin{table}[h]
\centering\small
\caption{All three detectors on the same alignments, sixteen taxa, five seeds,
counts of seeds significant at $0.05$. $\Phi_w$: window $50$ informative sites,
$199$ permutations, floor $1/200$. The quartet column is the certificate's margin
against a floor \emph{matched} to each population's own width, divergence and gap
load, decided by exact integer cross-product. Mean $\delta_q$ is calibrated by
parametric bootstrap under the same matched null, $99$ replicates, floor
$1/100$.}
\label{tab:headtohead}
\begin{tabular}{@{}llccc@{}}
\toprule
population & retic. & $\Phi_w$ & quartet vs matched & mean $\delta_q$ \\
\midrule
clonal, strict clock            & no  & $0/5$ & $\mathbf{0/5}$ & $0/5$ \\
clonal, \textbf{relaxed clock}  & no  & $\mathbf{0/5}$ & $\mathbf{0/5}$ & $0/5$ \\
one recombinant                 & yes & $5/5$ & $\mathbf{5/5}$ & $5/5$ \\
relaxed clock + one recombinant & yes & $5/5$ & $\mathbf{5/5}$ & $5/5$ \\
reciprocal mosaics              & yes & $5/5$ & $\mathbf{5/5}$ & $5/5$ \\
gap blocks                      & no  & \textbf{no measurement} & $\mathbf{0/5}$ & $0/5$ \\
saturated                       & no  & $\mathbf{0/5}$ & $\mathbf{0/5}$ & $0/5$ \\
\bottomrule
\end{tabular}
\end{table}

\paragraph{The result, stated plainly.} \textbf{All three detectors agree on all
seven populations} --- specific on the four without reticulation, detecting all
three with it at every seed --- except that $\Phi_w$ cannot be measured on the
gap-blocked data at all. \textbf{The incumbent is not beaten on power, and this
method's case cannot rest on power.} Edition~1 framed the missing comparison as
the validation we lacked; it is a three-way tie, and that is the more useful
finding.

Edition~2 of this table was wrong in our own favour's opposite direction, and the
correction is worth stating. It judged the certificate against the
\emph{strict-clock} population's floor while judging the competitors against
matched nulls, and so recorded the certificate as firing $5/5$ on the gap-blocked
and saturated populations --- neither of which contains reticulation. That was the
incommensurability this paper warns about in Section~\ref{sec:deletion},
committed a fourth time, against ourselves: the gap-blocked population has a
different denominator and the saturated one ten times the divergence. Given a
floor matched to its own width, divergence and gap load, the certificate is
specific on all four. Nothing about the method changed; the comparison was
corrected. The third column is the calibrated averaged statistic
(parametric bootstrap under a matched tree null), included because it is the
fairest possible competitor: same quadruples, same alignment, same matched-null
discipline, no permutation.

\paragraph{Two asymmetries, one each way.} The incumbent needs no recalibration
when divergence changes: it is specific on the saturated population at ten times
the divergence, where the certificate's raw margin against the strict-clock floor
fires at $5/5$. We are careful about what that $5/5$ is. The saturated
population's divergence and denominator are not the floor's, so the comparison is
not commensurable, and this is the same error --- in a fourth guise --- as
comparing margins across different denominators; the matched-floor construction
of Section~\ref{sec:real} exists precisely for it. But \emph{needing} a matched
null is a cost the incumbent does not pay, and that is a fair point against the
certificate in deployment. Conversely, on the gap-blocked population the
incumbent returns \textbf{no measurement at all}, having too few biallelic
gap-free informative columns once complete deletion is applied, where the
certificate still works.

\paragraph{What is left once power is off the table.} Two things, and we think
they are worth more than a power claim. First, a certificate is an
\textbf{exhibitable witness}: four sequences and six counts that a kernel
recomputes, so a reader re-verifies the finding without running our search or
trusting our description of it. Second, and specific to this paper's own
measurement: the null is a theorem rather than a permutation. The RDP5 manual
states that misalignment-induced false positives are detected with significant
$p$-values in the real alignment yet \emph{``will likely be undetectable in the
shuffled alignments''}~\cite{rdp5}. A site-order permutation null --- which is
exactly this incumbent's null --- is therefore blind by construction to the
artefact class Section~\ref{sec:aligner} measures. Neither point is established
by Table~\ref{tab:headtohead}; both are the actual case for the approach, and
stating them in place of a power claim is the honest position.

\section{Published breakpoints, and the resolution limit}\label{sec:breakpoint}

Section~\ref{sec:real}'s populations rest on two weak things: a literature
\emph{expectation} that a family recombines, and an alignment from the aligner of
Section~\ref{sec:aligner}, whose error had to be measured and gated on. Both can
be replaced, and doing so locates this method's resolution limit.

The Pango designation process records SARS-CoV-2 recombinants as \texttt{X*}
lineages with their parent lineages and, for many, an explicit numeric breakpoint
interval in reference coordinates, each traceable to a reviewed proposal; and
Nextstrain publishes an open release \emph{already aligned} by Nextclade to
MN908947, the frame those coordinates are quoted in. Two triads were built, 16
genomes each (five from each parent lineage, six from the recombinant), under
filters fixed in advance.

\begin{table}[h]
\centering\small
\caption{Published-breakpoint populations. The prediction, registered before the
test, is that the certificate's dominant pairing separates the two \emph{parent}
lineages with a recombinant on the other side --- a constraint on \emph{which}
four sequences are picked, which a margin cannot supply.}
\label{tab:breakpoint}
\begin{tabular}{@{}llccr@{}}
\toprule
recombinant & parents & breakpoint & prediction & margin \\
\midrule
XCP & EG.5.1.10 $\times$ JD.1.1 & 14190--14633 & \textbf{held}, 6/10 & $2/28920$ \\
XCM & XBB.2.3 $\times$ DV.7.1 & 22017--22031 & \textbf{failed}, 0/10 & $4/28511$ \\
\bottomrule
\end{tabular}
\end{table}

On XCP the method picked the two true parents as the far pair in six of its ten
strongest certificates. On XCM every certificate instead paired a parent with a
recombinant on each side. One of two is a split result and is reported as one.

\paragraph{The finding is a resolution limit.} Both margins are at the edge of
what the statistic resolves. XCP's population has mean $p$-distance $0.00057$ and
XCM's $0.00127$, implying roughly $17$ and $38$ differing sites across $29\,903$
columns; the margins are $2$ and $4$ \emph{sites}. A verdict cannot rest on that
however the floor comparison comes out --- and XCP's floor comparison is itself
degenerate, the matched null having produced no violating quadruple in any
replicate, so ``exceeds the floor'' reduces to ``the margin is positive''. Both
the degeneracy and the site counts are now computed into the output rather than
left to prose.

So, stated against our own interest: \textbf{SARS-CoV-2 recombinants between
recently-diverged lineages sit below this method's resolution.} And there is a
tension worth naming because it will recur --- \textbf{the datasets with the best
ground truth are the least divergent}, since recombinants are named between
lineages close enough to be confidently assigned. Section~\ref{sec:real}'s
populations, at mean $p$-distance $0.0088$ and $0.11$, suit the statistic far
better and are far worse supplied with truth.

One signal survives, labelled as suggestive rather than evidence: on XCP the
\emph{pairing} was right in 6 of 10 certificates, in a regime where the magnitude
means nothing. The method pointed at the correct two parents.

\paragraph{A note on the trust surface.} Using a published alignment removes
\emph{our} aligner and substitutes Nextclade's, which is reference-anchored too
--- the same construction and so the same class of error Section~\ref{sec:aligner}
measures. Provenance improves; measurement coverage does not, since we have not
measured their error. The admissibility gate calibrated for ours does not
transfer and is not applied.

\section{The maximum statistic's defect, measured across taxon counts}
\label{sec:scale}

Section~\ref{sec:maxmean} concluded that the certificate's weakness is not power
but incomparability across taxon counts, and argued it from two data points. This
paper has twice recorded the cost of reasoning from too few points, so the defect
is measured.

That required making the measurement affordable. Sixty-four taxa is $635\,376$
quadruples and the byte-at-a-time counting visits $4000$ columns per quadruple.
Packing each sequence into three bitvectors over columns --- the base's two bits,
plus a mask of definite bases --- turns a quadruple's seven counts into one pass
over $L/64$ words: intersect four masks for the common columns, then count
differing bits for each of the six pairs inside it. Measured $21.6\times$ faster
at thirty-two taxa with identical exact margins. It is a second implementation,
not a replacement: the byte path remains the specification, agreement is a
property test across widths $1$, $7$, $63$, $64$, $65$, $130$ and $257$ plus
gap-blocked data, and the sweep re-checks it on every run.

\begin{table}[h]
\centering\small
\caption{Exhaustive over every quadruple, three seeds, margins over $4000$. The
``above floor'' column is seed-paired and decided by exact integer
cross-product.}
\label{tab:scale}
\begin{tabular}{@{}rrllcl@{}}
\toprule
taxa & quadruples & clonal floor & one recombinant & above floor & mosaics \\
\midrule
16 & $1\,820$ & 7, 5, 6 & 60, 52, 53 & $3/3$ & 136, 112, 124 \\
32 & $35\,960$ & 10, 9, 10 & 60, 67, 98 & $3/3$ & 148, 154, 108 \\
64 & $635\,376$ & 17, 16, 14 & 89, 82, 83 & $3/3$ & 156, 148, 166 \\
\bottomrule
\end{tabular}
\end{table}

The floor roughly doubles from sixteen to sixty-four taxa while the signal stays
flat, so the worst-seed separation of mosaics over the floor halves:
$19.4\times$, $10.8\times$, $9.2\times$. Two readings, both needed. Detection at
a \emph{fixed} taxon count survives comfortably. But ``twice the floor'' means
something different at sixteen taxa than at sixty-four, so a threshold
calibrated at one size does not transfer --- which is why
Section~\ref{sec:real}'s verdicts use a floor matched to each population, and
why edition~2's raw-margin comparison was the wrong one.

One figure looks alarming and is not: at sixty-four taxa $82\%$ of all
quadruples violate the four-point condition on a clonal, strict-clock
population. Finite sequences are never exactly additive, which is precisely why
the \emph{margin} and not the \emph{count} carries the inference --- and why the
proportion-violating variant of Section~\ref{sec:maxmean} separated by only
$1.22\times$.

\section{A topology experiment, and the grounds for closing it}
\label{sec:topology}

The construction this programme began from reads an unfilled $1$-cycle in the
Vietoris--Rips complex of a genetic-distance cloud as reticulation, on the same
warrant as our null: a tree metric admits no unfilled cycle. A search of the
literature found no peer-reviewed independent calibration of that method's
specificity. Our grid has the ground truth to supply one, since the saturated
dataset differs from the baseline only in having ten times the substitution rate.

At sixteen taxa the experiment is \textbf{inconclusive} and its output says so
(\texttt{"this\_experiment\_is\_interpretable": false}): the matched null produced
no cycle in $495$ of $495$ replicates on five of seven datasets, so the $p$-value
takes only two values and the test reduces to \emph{did a cycle appear at all}.
The cause is diagnosable --- at that size the complex is too small for $H_1$ to
populate --- and the fix is more taxa, which \S\ref{sec:scale} made affordable.

Re-run at sixty-four taxa, the answer arrives. Two nulls become
non-degenerate, and they are the two datasets containing \emph{no reticulation}:
the statistic fires at $3/3$ on gap-blocked data and at $3/3$ on saturated data,
while a genuine single recombinant produces \textbf{no cycles at all} and is
missed at $0/3$.

Five hypotheses were then written down for that failure, three selected on
cost-to-decide, and the criterion fixed in advance: a variant succeeds if the
reciprocal-mosaic population is significant at every seed and both
no-reticulation populations at none.

\begin{table}[h]
\centering\footnotesize
\caption{Filtration funnel, sixty-four taxa, three seeds, 49 bootstrap
replicates. No variant meets the criterion.}
\label{tab:funnel}
\begin{tabular}{@{}llcccc@{}}
\toprule
variant & hypothesis & mosaics & one recomb. & gaps & saturated \\
\midrule
baseline, total & --- & $3/3$ & $\mathbf{0/3}$ & $2/3$ & $3/3$ \\
F1 Jukes--Cantor & saturation & $3/3$ & $\mathbf{0/3}$ & $2/3$ & $3/3$ \\
F2 global complete deletion & denominators & $3/3$ & $\mathbf{0/3}$ & $\mathbf{0/3}$ & $3/3$ \\
F3 longest bar & the summary & $3/3$ & $\mathbf{0/3}$ & $\mathbf{0/3}$ & $3/3$ \\
\bottomrule
\end{tabular}
\end{table}

\textbf{F1 is refuted}: correcting for saturation leaves that cell at $3/3$ and
\emph{raises} its statistic, $0.00867 \to 0.01421$. \textbf{F2 is confirmed}, and
confirms the construction argument below: global complete deletion takes
gap-blocked data from $13.3$ cycles and $2/3$ to $0.0$ cycles and $0/3$, so the
gap false positive was the missing-data confound of Section~\ref{sec:deletion}
reappearing as topology. F3 removes it too, by distinguishing many weak cycles
from few strong ones. But \textbf{the miss is not a filtration artefact} --- zero
cycles on a single recombinant in all four variants --- and the certificate
detects that same population at $3/3$ with roughly fivefold margins
(Table~\ref{tab:scale}).

\paragraph{Decision: the line is closed.} Not because topology is the wrong
idea; the warrant, \emph{a tree metric admits no unfilled cycle}, is the same
fact this paper's null rests on. What is closed is the
global-distance-matrix implementation, for two reasons the funnel established
rather than assumed. It \textbf{cannot use per-subset complete deletion}, because
a complex is built from all pairs at once and there is no ``within this subset''
to restrict to; F2 shows the available substitute works but costs every column
any taxon is missing. And it \textbf{misses the signal the certificate catches}.
So this programme's topological content is the certificate itself: the same
theorem, evaluated combinatorially per quadruple, which is exactly what makes
complete deletion within the subset available at all.

\paragraph{Limits.} Three seeds, one generative model, and the saturated cell's
null is degenerate in every variant, so that column carries no calibration and
its $3/3$ is the binary above. The two conclusions the decision rests on --- F1's
refutation and the single-recombinant miss --- do not depend on it.

\begin{remark}[The lesson both of these share]
Two experiments in this edition produced a confident-looking count from a null
that could not vary. A null that is identically zero makes every positive
observation significant at the floor, which is indistinguishable from a powerful
test until you look at the null. Both now compute a degeneracy flag, and no
calibrated statistic in this work should have its $p$-value read before that flag
is checked.
\end{remark}

\section{Real populations}\label{sec:real}

Five viral taxa were drawn from a public corpus of $58\,046$ viral genomes with
taxonomy~\cite{vhg} under filters fixed in advance: within $2\%$ of the taxon's
median length, at most $0.1\%$ ambiguous residues, exact duplicates removed,
sixteen sequences per population, selected either in accession order or by greedy
$6$-mer diversity. Each taxon carries a \emph{literature} expectation about
whether its genealogy is tree-like, written before any margin was computed, and
the test is free to contradict it. Ten populations result.

Seven fail the admissibility gate of Section~\ref{sec:aligner} and are reported
blocked with their counts rather than interpreted. This is expensive on purpose:
one discarded population has a margin of $30/1784$ and another $13/5082$. Note
also which selection strategy survives --- diversity maximisation is exactly what
breaks reference-anchored alignment. Each of the three that pass is aligned twice,
on the two most central sequences, and each margin is compared against a floor
from a simulated clonal relaxed-clock tree at \emph{matched} alignment width,
\emph{matched} mean $p$-distance (tuned by bisection) and \emph{matched} gap
load, because a margin alone is not a statistic and a floor from the wrong
divergence measures the floor.

\begin{table}[h]
\centering\small
\caption{The three admissible populations. References agree on the verdict in all
three.}
\label{tab:real}
\begin{tabular}{@{}llrrl@{}}
\toprule
population & expectation & margin & matched floor & verdict \\
\midrule
Hepatitis B, accession order & reticulate & $9/3182$ & $2/3182$ & above $2\times$ \\
Coxsackievirus A6, diversity & reticulate & $329/7312$ & $268/7390$ & above $1\times$ only \\
JC polyomavirus, accession order & tree-like & $4/5108$ & $24/5121$ & \textbf{below} \\
\bottomrule
\end{tabular}
\end{table}

Three populations, three agreements with the prior. The JC polyomavirus row is
the one that could most easily have embarrassed the method and did not: its
margin is a sixth of the floor a tree produces at matched width and divergence.
Both populations clearing their floor carry a Lean file embedding the four
aligned rows, so the kernel recomputes the six pair counts and the common column
count itself: accepted at $3182$ columns in $52.7$~s and at $7434$ columns in
$117.6$~s, with both corrupted twins rejected.

Two cautions belong beside Table~\ref{tab:real} rather than in a footnote. Three
is not a validation. And the Coxsackievirus row is weak and at high divergence
(mean $p = 0.11$): it clears $1\times$ and not $2\times$, and the noise floor
itself grows with divergence, so it should be read as not inconsistent with the
prior rather than as a detection. A further caution concerns the prior itself:
one of our two negative-control families, the papillomaviruses, has documented
within-type recombination~\cite{murahwa}, so that prior is the weakest in the set
and we mark it as such. A clonal bacterial lineage would be the stronger negative
control.

\section{The Gromov hinge of Buneman's converse}\label{sec:gromov}

Theorem~\ref{thm:buneman} is cited above and not proved. The constructive proof
in the literature does not attack the tree directly: it factors through the
\emph{Gromov product} at a chosen root, reducing the four-point condition on $d$
to an ultrametric-style condition on a derived function, and only then builds an
equidistant tree by a cherry-picking recursion. The reduction is the hinge, and
unlike the tree construction it is pure algebra.

\begin{definition}
At a root $r$, write
$g_r(x,y) = \tfrac{1}{2}\big(d(x,y) - d(r,x) - d(r,y)\big)$. This is a
dissimilarity, not a metric: it takes negative values. Say $f$ satisfies the
\emph{strong triangle inequality} when $f(x,z) \le \max(f(x,y), f(y,z))$ for all
$x,y,z$ --- inequality \eqref{eq:ultra} with no positivity required.
\end{definition}

\begin{theorem}[The hinge; proved, not cited]\label{thm:gromov}
A symmetric $d$ satisfies the four-point condition if and only if $g_r$ satisfies
the strong triangle inequality for every root $r$.
\end{theorem}

\begin{proof}
Expanding $g_r(x,z) \le g_r(x,y)$ gives
$d(x,z) + d(r,y) \le d(x,y) + d(r,z)$, and $g_r(x,z) \le g_r(y,z)$ gives
$d(x,z) + d(r,y) \le d(y,z) + d(r,x)$. So the right-hand side of the theorem says
exactly
\[
  d(x,z) + d(r,y) \;\le\;
  \max\!\big(d(x,y)+d(r,z),\; d(y,z)+d(r,x)\big)
\]
for all $r,x,y,z$, which is \eqref{eq:fp} instantiated at
$(a,b,c,e) = (x,z,r,y)$ together with the symmetry of $d$. Reading the
implication backwards requires the root $c$ with the remaining three in the order
$(a,e,b)$.
\end{proof}

\begin{lstlisting}[caption={The Lean statement. \texttt{\#print axioms} gives
\texttt{[propext, Classical.choice, Quot.sound]}.},label={lst:gromov}]
theorem fourPoint_iff_gromov_strongTriangle {d : α → α → ℚ}
    (hsym : ∀ x y, d x y = d y x) :
    IsFourPoint d ↔ ∀ r, StrongTriangle (gromov d r)
\end{lstlisting}

Two details are worth recording. The condition is named \texttt{StrongTriangle}
and not \texttt{IsUltrametric}, because the Gromov product takes negative values
and calling it an ultrametric would import a positivity hypothesis it does not
satisfy; the literature's phrase ``the Gromov product is an ultrametric'' means
only the max inequality. And the kernel earned its place: our first attempt
instantiated the reverse direction at the wrong permutation, \texttt{linarith}
failed on two goals, and the axiom gate reported \texttt{sorryAx} rather than
letting a plausible-looking proof through.

What remains cited is the tree construction --- the equidistant cherry-picking
recursion and the re-rooting that returns a tree metric. The chain this paper can
claim is therefore: certificate $\Rightarrow$ not four-point $\Rightarrow$
(proved) some root's Gromov product fails the strong triangle inequality
$\Rightarrow$ (cited) not realised by a weighted tree. We searched Mathlib, the
Isabelle Archive of Formal Proofs and the Coq ecosystem and found ultrametric
spaces but no tree metrics, no splits and no four-point condition, so completing
the construction appears to be first of its kind. That search was across indexes
rather than exhaustive.

\section{Limitations, and what would make this wrong}\label{sec:limits}

\begin{itemize}
\item \textbf{Buneman's converse is half-cited.} Theorem~\ref{thm:gromov} closes
the first step; the construction is not formalised. So ``satisfies \eqref{eq:fp}
everywhere'' is not yet formally ``is a tree'' in this development.
\item \textbf{A certificate does not identify an event.} Satisfying \eqref{eq:fp}
does not certify the \emph{absence} of reticulation, and a violation does not
pin down which tree the data departs from; there exist networks whose distances
are representable by a tree other than their own underlying
tree~\cite{francissteel}. The honest report is \emph{that} the data are not
tree-like and \emph{where}, not an event history.
\item \textbf{The margin's null distribution is not derived.} Finite sequences
are never exactly ultrametric, so certificates exist under the null with small
margins; we calibrate by matched mechanistic simulation rather than by a
distribution. We found no analytic sampling distribution of the four-point margin
in the literature either, only high-probability sequence-length bounds and
parametric bootstrap, so this is the field's position and not only ours.
\item \textbf{Three admissible real populations.} Too few to call the method
validated, and the aligner that made them admissible is reference-anchored.
\item \textbf{The head-to-head is a tie, not a win}
(Section~\ref{sec:headtohead}), and it covers one incumbent. The comparison
against the quartet-averaged treelikeness statistics \emph{used as calibrated
hypothesis tests} rather than as scores has still not been run, nor against the
clock-assuming tests, which would fail Section~\ref{sec:refute}'s control but
are what most practitioners actually use.
\item \textbf{The statistic is an extreme value.} Section~\ref{sec:maxmean}
prices this at fixed taxon count only.
\end{itemize}

What would make the programme wrong, rather than merely incomplete: if
violations of \eqref{eq:fp} at the scale we report turned out to be dominated by
some departure from tree-likeness other than reticulation --- population
structure, recurrent mutation, or selection acting locally --- then the
certificates would remain valid and the biological reading of them would not. We
have controlled for rate variation and for missing data, and measured the
alignment contribution. We have not controlled for the rest.

\section{Reproducibility}\label{sec:repro}

Every number in this paper is produced by a named command in the repository
\texttt{SocrateAI-Scientific-Genetic} at release \texttt{v1.8.0}, and the raw
output is deposited with it.

\begin{center}\small
\begin{tabularx}{\textwidth}{@{}lX@{}}
\toprule
command & what it produces \\
\midrule
\texttt{reticulation experiment} & Table~\ref{tab:refute},
Section~\ref{sec:deletion}, Table~\ref{tab:maxmean} \\
\texttt{reticulation align-error} & Table~\ref{tab:ladder} and the admissibility
threshold it measures \\
\texttt{reticulation real-data} & Table~\ref{tab:real} and the kernel verdicts \\
\texttt{reticulation saturation} & the length-scaling diagnostic behind the
divergence caution \\
\texttt{reticulation experiment} \texttt{-{}-phi-window}, \texttt{-{}-phi-permutations}
& Table~\ref{tab:headtohead}; defaults $50$ and $199$ \\
\texttt{reticulation scale} & Table~\ref{tab:scale} and the $21.6\times$ check \\
\texttt{reticulation filtration} & Table~\ref{tab:funnel} \\
\texttt{reticulation breakpoint} & Table~\ref{tab:breakpoint} \\
\texttt{scripts/reticulata\_hf\_populations.py} & the ten populations, with a
SHA-256 of the source file \\
\texttt{lake build ANSE.GeneticaReticulataTree} & the theorems, with
\texttt{\#print axioms} on each \\
\bottomrule
\end{tabularx}
\end{center}

All are deterministic given their seeds. The population extraction is
deterministic given the source corpus, whose digest, filters and every selected
accession are recorded in \texttt{provenance.json}. Verification at the release:
$98$ library and $2$ integration tests in Rust; Lean $18/18$ and $14/14$
declarations with zero \texttt{sorry} and no unwhitelisted axioms; $33$ Python
tests; an anti-stub audit over $227$ files.

\paragraph{Status of this manuscript.} Not peer reviewed. The repository
previously described work as ``5-round Opus peer reviewed'' on the basis of a
function that made no model call; that description was withdrawn, and no
equivalent claim is made here. Reported results are the output of the commands
above, including the three places where that output contradicted what we had
written.

\begin{thebibliography}{99}

\bibitem{buneman} Buneman P (1971/1974). The recovery of trees from measures of
dissimilarity. In \emph{Mathematics in the Archaeological and Historical
Sciences}, Edinburgh University Press, 387--395; and: A note on the metric
properties of trees. \emph{J. Combinatorial Theory B} \textbf{17}:48--50.
[Attribution is contested between these two in the secondary literature; neither
was retrievable as a primary record here, and no DOI is asserted.]

\bibitem{semplesteel} Semple C, Steel M (2003). \emph{Phylogenetics}. Oxford
Lecture Series in Mathematics and its Applications 24, Oxford University Press.
[The statement of Theorem~\ref{thm:buneman} used here was taken from S.~Roch's
MATH~833 notes, which follow this book's Chapter~7; we do not assert a theorem
number.]

\bibitem{francissteel} Francis AR, Steel M (2014). Tree-like reticulation
networks --- when do tree-like distances also support reticulate evolution?
\emph{arXiv:1405.2965}.

\bibitem{gibrat} Gibrat J-F (2018). A short note on dynamic programming in a
band. \emph{BMC Bioinformatics} \textbf{19}:226.
\href{https://doi.org/10.1186/s12859-018-2228-9}{doi:10.1186/s12859-018-2228-9}

\bibitem{flouri} Flouri T, Kobert K, Rognes T, Stamatakis A. Are all global
alignment algorithms and implementations correct? \emph{bioRxiv} preprint.
\href{https://doi.org/10.1101/031500}{doi:10.1101/031500}

\bibitem{bruen} Bruen TC, Philippe H, Bryant D (2006). A simple and robust
statistical test for detecting the presence of recombination. \emph{Genetics}
\textbf{172}(4):2665--2681.
\href{https://doi.org/10.1534/genetics.105.048975}{doi:10.1534/genetics.105.048975}

\bibitem{deltaplots} Holland BR, Huber KT, Dress A, Moulton V (2002). $\delta$
plots: a tool for analyzing phylogenetic distance data. \emph{Mol. Biol. Evol.}
\textbf{19}(12):2051--2059.
\href{https://doi.org/10.1093/oxfordjournals.molbev.a004030}{doi:10.1093/oxfordjournals.molbev.a004030}

\bibitem{frankelane} Frankel LE, An\'e C (2023). Summary tests of introgression
are highly sensitive to rate variation across lineages. \emph{Systematic Biology}
\textbf{72}(6):1357--1369.
\href{https://doi.org/10.1093/sysbio/syad056}{doi:10.1093/sysbio/syad056}

\bibitem{murahwa} Murahwa AT, \emph{et al.}\ (2020). Recombination between
high-risk human papillomaviruses and non-human primate papillomaviruses.
\emph{J. Mol. Evol.} \textbf{88}:453--462.
\href{https://doi.org/10.1007/s00239-020-09946-0}{doi:10.1007/s00239-020-09946-0}

\bibitem{rdp5} Martin DP, \emph{et al.} \emph{RDP5 Instruction Manual}.
\url{https://web.cbio.uct.ac.za/~darren/RDP5Manual.pdf} [Quotations are verbatim
from §3.1.2, §3.1.3, §10.2 and §12.1 of the manual as retrieved.]

\bibitem{jaya} Jaya FR, Brito BP, Darling AE (2023). Evaluation of
recombination detection methods for viral sequencing. \emph{Virus Evolution}
\textbf{9}(2):vead066.
\href{https://doi.org/10.1093/ve/vead066}{doi:10.1093/ve/vead066}
[Cited here for the fact that the benchmark does not vary alignment quality.]

\bibitem{vhg} hiyata (2025). \emph{Virus-Host-Genomes}. Hugging Face dataset,
$58\,046$ viral genomes with taxonomy, MIT licence.
\url{https://huggingface.co/datasets/hiyata/Virus-Host-Genomes}

\bibitem{chan} Chan JM, Carlsson G, Rabadan R (2013). Topology of viral
evolution. \emph{PNAS} \textbf{110}(46):18566--18571.
\href{https://doi.org/10.1073/pnas.1313480110}{doi:10.1073/pnas.1313480110}
[The construction this work follows: persistence on a population under a
genetic-distance metric.]

\end{thebibliography}

\end{document}
